
\newpage  % 开始新的一页放置正文内容
\section{\huge \textcolor{blue}{几何视角下的微分方程}}
\begin{center}
\huge GEOM VIEW OF ODS'S
\end{center}

\begin{center}
\Large Analyze(解析) \hspace{3cm} Geometrically(几何)
$y' = (x, y) \xleftrightarrow{\hspace{3cm}} \text{方向场 (direction field)}$
$y_1 :solution \xleftrightarrow{\hspace{2cm}} \text{积分曲线 (Integral curves)}$
\end{center}
\vspace{1cm}
\subsection{\LARGE 方向场}
\large direction field 

\begin{figure}[htbp]
    \centering
\includegraphics[width=1\textwidth]{C:/Users/37/Desktop/Notes on Common Calculus Equations/Notes and pictures/cwjffc1.png}
    \caption{方向场示意图}
    \label{fig:direction-field}
\end{figure}

\newpage
\sethlcolor{yellow}
\text{有如下关系}\\
$y'(x) \text{ is solution to } y'=f(x,y)\Longleftrightarrow \text{graph of y'(x) is an Int curves}$
$y_1'(x)=f(x,y_1(x))\hspace{2.2cm} \Longleftrightarrow {\hspace{2cm}} \text{shape of } y_1'(x)$\\[1cm]
简单来说，即在解析法中写出微分方程等价于几何法中画出方向场,而解析法中求出微分方程的解等价于几何法中作出积分曲线图像.
\vspace{1cm}
\subsection{\LARGE DRAWINGQ DIRECTION FIELD}

\subsubsection{\large 计算机}

\begin{enumerate}  %使用 enumerate 环境实现编号列表对齐
\item Pick(x,y)   等距
\item $f(x,y) \longrightarrow \text{求值}$
\item 画 $f(x,y)$ 线素(可以理解为斜率)
\end{enumerate}

\subsubsection{\large 人}
\begin{enumerate}  %使用 enumerate 环境实现编号列表对齐
\item 取固定斜率 C
\item 满足方程 $f(x,y)=C$  ，画出方程\\
level curve(等值线) or isocline(等斜线)\\
在微积分中称此方程所描绘曲线为等值线,而在ODS中称为等斜线.
\item slope C\\
画出等斜线(一般用虚线表示，因为等斜线并非为方程解)，并在等斜线上画出一定量的值为c的线素

\newpage
\ 下用一例题演示作图流程

\end{enumerate}

\begin{figure}[htbp]
    \centering
\includegraphics[width=0.7\textwidth]{C:/Users/37/Desktop/Notes on Common Calculus Equations/Notes and pictures/cwffc2.png}
    \caption{例题1}
    \label{fig:eg1}
\end{figure}

\ 由方向场画得积分曲线图，显然得一同心圆族，即可得积分曲线方程为
{\Large
\begin{equation}
x^2 + y^2 = C^2
\end{equation}
}
\ 为更一般得了解几何绘图，继续给出以下例子

\begin{figure}[htbp]
    \centering
\includegraphics[width=0.8\textwidth]{C:/Users/37/Desktop/Notes on Common Calculus Equations/Notes and pictures/cwffc3.png}
    \caption{例题2}
    \label{fig:eg2}
\end{figure}

\newpage
\ 观察例题2中夹在 C=0 和 C=2 等斜线之间的区域，可以发现积分曲线在此区域由方向场的引导逐渐靠近 C=1 的等斜线，且无法逃逸.\\
%isoclines(等斜线)
对2条积分曲线不相交现象, 进行探讨并给出如下原理

\subsection{\Large 积分曲线作图原理}
\subsubsection{\large 积分曲线作图原理 1}
%使用 enumerate 环境实现编号列表对齐
\begin{enumerate}
\Large
\item 两条积分曲线不相交
\item 存在与唯一性定理(不相切)
\end{enumerate}

对1: 如果相交则会在交点产生2个线素, 产生矛盾\\
\hspace*{2em}对2: 相对于1更不显然, 表述为2条积分曲线不仅不能相交, 甚至连相切
都无法做到，为 $ y'=(x,y)  $  有且仅有一个解.  有, 即存在性,仅有,即唯一性.\\
其中y为连续函数保证存在性成立, 
那要如何保证唯一性的假设? \ 即y的偏导必须连续(均在$(x_0,y_0)$附近成立)\\[0.5cm]
在了解曲线积分作图原理后, 再回头看例题2\\
\hspace*{2em}积分曲线由方向场的引导逐渐靠近 C=1 时得到的等斜线，曲线无法逃逸出直线 y=x 附近
且直线两侧的积分曲线不能交叉、接触，则我们可以用几何方法猜测 y=x 为微分方程的一个解析解.\\[0.5cm]

\newpage
 由上述原理再次思考前例题后，我们再看下例题
\begin{figure}[htbp]
    \centering
\includegraphics[width=1\textwidth]{C:/Users/37/Desktop/Notes on Common Calculus Equations/Notes and pictures/cwffc4.png}
    \caption{例题3}
    \label{fig:eg3}
\end{figure}

\ 现观察例题3, 解在除了x=0 直线外的平面上，下观察此直线,
P点为方程所生产的所有直线交点, 即在点P上唯一性不成立，再看直线 x=0 外其他点，
无直线经过，即存在性不成立.\\
那么我们发现，此解法不满足原理，产生矛盾，那问题会在哪?\\
我们对方程一开始的变形做出更改，以至于符合标准形式:
{\Large
\begin{equation}
\frac{dy}{dx} = \frac{1-y}{x}
\end{equation}
}
\ 此时我们即发现方程在 x=0 处无定义, 即存在及唯一性定理不保证在y轴(x=0)上成立，则说明上述并不矛盾.\\
同时此例题也提醒我们在利用原理时要注意方程是否满足标准形式，再去分析.
